\section*{अध्याय \ref{ch:b2:pt-paths} --- तापदाबमिति और P--T--t पथ}

\begin{solution}{exo:b2:pt-paths:1}
(क) A क्लोराइट मंडल, B बायोटाइट मंडल, C गार्नेट मंडल, D स्टॉरोलाइट मंडल, E कायनाइट मंडल,
F सिलिमेनाइट मंडल। (ख) पाँच: बायोटाइट, गार्नेट, स्टॉरोलाइट, कायनाइट और सिलिमेनाइट की
समश्रेणी-रेखाएँ, हर एक दो क्रमागत नमूनों के बीच, और श्रेणी A से F तक बढ़ती है। (ग) समश्रेणी-रेखा
केवल यह बताती है कि कोई अभिक्रिया हुई; क्षेत्र-प्रवणता को दाब और तापमान चाहिए, जो
अंशांकित अभिक्रियाओं से या हर चट्टान के शिखर-समुच्चय की तापदाबमिति से
आते हैं।
\end{solution}

\begin{solution}{exo:b2:pt-paths:2}
(क) $(X_{\text{Mg}}/X_{\text{Fe}})_{\text{गार्नेट}} = 0.10/0.80 = 0.125$,
$(X_{\text{Mg}}/X_{\text{Fe}})_{\text{बायोटाइट}} = 0.45/0.50 = 0.900$; $K_D = 0.1389$, $\ln K_D =
-1.9741$। (ख) $3R = 24.942$; अब हर का मान है $20.19 + 24.942\times1.9741 = 69.428$;
\qty{5}{kbar} पर $T = 54185/69.428 = \qty{780.5}{K} = \qty{507}{\celsius}$;
\qty{7}{kbar} पर $T = (52630 + 2177)/69.428 = \qty{789.4}{K} = \qty{516}{\celsius}$।
(ग) दो किलोबार उसे केवल \qty{9}{K} बदलते हैं: एक तापमापी।
\end{solution}

\begin{solution}{exo:b2:pt-paths:3}
(क) $z = P/(\bar\rho g) = 6.0\times10^8/(2800\times9.81) = \qty{21.8}{km}$।
(ख) $(580 - 10)/21.8 = \qty{26}{\celsius\per\km}$। (ग) \qty{580}{\celsius} पर रेखा
$3.67 + 0.8\times2.10 = \qty{5.35}{kbar}$ पर है; \qty{6.0}{kbar} पर चट्टान
उच्च-दाब पक्ष पर है: कायनाइट।
\end{solution}

\begin{solution}{exo:b2:pt-paths:4}
(क) अभिकारक $2\times140.693 + 3\times18.015 = \qty{335.431}{g}$; उत्पाद $277.112 +
58.320 = \qty{335.432}{g}$: पूर्णांकन की सीमा तक संतुलित। (ख) $2\times43.65 =
\qty{87.30}{cm^3}$ बन जाते हैं $107.50 + 24.63 = \qty{132.13}{cm^3}$: $+\qty{51.4}{\percent}$।
(ग) ओलिवीन $281.386/87.30 = \qty{3.223}{g.cm^{-3}}$; उत्पाद $335.432/132.13 =
\qty{2.539}{g.cm^{-3}}$। (घ) चट्टान अपने आयतन के आधे जितनी फूल जाती है और हल्की हो जाती है; वह
अपने खनिजों में जल भी संचित करती है।
\end{solution}

\begin{solution}{exo:b2:pt-paths:5}
(क) माध्य $-1.206$, \qty{1072.15}{K} पर, और $-1.782$, \qty{823.15}{K} पर। $1/T =
9.3271\times10^{-4}$ और $1.21485\times10^{-3}$ के साथ, $a = (-1.206 + 1.782)/(9.3271\times10^{-4}
- 1.21485\times10^{-3}) = -2041.5$ K और $b = -1.206 + 2041.5/1072.15 = 0.6982$।
(ख) $\Delta_r H + 2070\times0.311 = 2041.5\times24.942 = \qty{50920}{J}$, अतः $\Delta_r H =
\qty{50276}{J}$; $\Delta_r S = 0.6982\times24.942 = \qty{17.41}{J.K^{-1}}$। (ग)
\qty{5}{kbar} पर: $K_D = 0.18$ के लिए 588.1, \qty{587.5}{\celsius} की तुलना में; $0.10$ के लिए 419,
\qty{425}{\celsius} की तुलना में; $0.30$ के लिए 819, \qty{806}{\celsius} की तुलना में। प्रयोगों के परास
के भीतर दोनों अंशांकन सहमत हैं, यद्यपि उनके स्थिरांक
\qty{2.4}{kJ} और \qty{2.8}{J.K^{-1}} से भिन्न हैं; $\Delta_r H$ और $\Delta_r S$ की त्रुटियाँ
वहाँ एक-दूसरे की क्षतिपूर्ति करती हैं और उसके बाहर बढ़ती हैं। किसी अंशांकन का उपयोग उसके प्रयोगों की
दशाओं से बहुत बाहर नहीं किया जाना चाहिए।
\end{solution}

\begin{solution}{exo:b2:pt-paths:6}
(क) विनिमय: $(113.12 + 154.30) - (115.32 + 149.65) = \qty{+2.45}{cm^3} =
\qty{0.245}{J.bar^{-1}}$, 0.311 की तुलना में। GASP: $125.28 + 2\times44.15 + 22.69 -
3\times100.79 = \qty{-66.10}{cm^3} = \qty{-6.61}{J.bar^{-1}}$, $-6.21$ की तुलना में, जो
\qty{800}{\celsius} और \qty{10}{kbar} पर अंशांकन का मान है, जहाँ ऊष्मीय प्रसार और
संपीडन आयतनों को बदलते हैं। दोनों स्थितियों में GASP परिवर्तन विनिमय वाले का बीस से तीस गुना
है। (ख) $\partial T/\partial P = 0.311/62.96 = \qty{4.9e-3}{K/bar}$:
प्रति किलोबार \qty{4.9}{K}। (ग) रेखा की ढलान $22.80 + 8.314\times3\ln 0.20/6.21
= 22.80 - 6.46 = \qty{16.3}{bar/K}$ है; \qty{17}{K} दाब को
\qty{280}{bar} खिसकाते हैं।
\end{solution}

\begin{solution}{exo:b2:pt-paths:7}
(क) $P_0 = 22.80\times600 - 1093 = \qty{12587}{bar}$; $3\ln(0.08/0.30) = -3.9653$; $RT =
8.314\times873.15 = \qty{7259.4}{J}$; $P = 12587 + 7259.4\times(-3.9653)/6.21 = 12587 -
4635 = \qty{7952}{bar}$। (ख) $3\ln(0.08/0.35) = -4.4277$, $P = 12587 - 5176 =
\qty{7411}{bar}$: \qty{0.54}{kbar} कम, इसलिए प्लैजियोक्लेज़ का विश्लेषण सावधानी से करना चाहिए।
(ग) $7952\times10^5/(2800\times9.81) = \qty{29.0}{km}$।
\end{solution}

\begin{solution}{exo:b2:pt-paths:8}
(क) $z = P/(\bar\rho g)$: \qty{118.9}{km}, \qty{34.0}{km} और \qty{15.3}{km}।
(ख) $(118.9 - 34.0)/2.2 = \qty{38.6}{km/Myr} = \qty{3.9}{cm.yr^{-1}}$ और $(34.0 -
15.3)/1.1 = \qty{17.0}{km/Myr} = \qty{1.7}{cm.yr^{-1}}$, प्रकाशित 3.4 और 1.6 से थोड़ा
ऊपर: उनका दाब से गहराई में रूपांतरण गहराई के थोड़े छोटे
अंतर देता है। (ग) $\sigma = \sqrt{0.9^2 + 0.9^2} = \qty{1.3}{Myr}$, अतः $2.2 \pm 1.3$
मिलियन वर्ष और $84.9/3.47 = 24$ तथा $84.9/0.93 = \qty{91}{km/Myr}$ के बीच की दर,
यानी 2.4 से \qty{9.1}{cm.yr^{-1}}: हर स्थिति में तेज़। (घ) सबसे कम दर भी मॉडल के \qty{0.5}{mm.yr^{-1}}
की लगभग पचास गुनी है: अकेला अपरदन यह नहीं कर सकता।
\end{solution}

\begin{solution}{exo:b2:pt-paths:9}
(क) $\partial T/\partial\ln K_D = 3RT/(\Delta_r S - 3R\ln K_D) = 24.942\times860.6/62.96 =
\qty{341}{K}$; $\partial T/\partial P = 0.311/62.96 = \qty{4.9}{K}$ प्रति किलोबार।
(ख) $\sigma_T = \sqrt{(341\times0.05)^2 + 4.9^2} = \sqrt{17.0^2 + 4.9^2} =
\qty{17.7}{K}$। (ग) $K_D$ पर विश्लेषणात्मक त्रुटि प्रभावी है; वह अकेले सबसे अच्छे अंशांकनों के
\qty{25}{\celsius} का दो-तिहाई है, जिनमें अंशांकन की अपनी त्रुटियाँ भी
शामिल हैं।
\end{solution}

\begin{solution}{exo:b2:pt-paths:10}
(क) गार्नेट $0.10/0.75 = 0.1333$, बायोटाइट $0.42/0.52 = 0.8077$, $K_D = 0.1650$, $\ln K_D =
-1.8018$, हर $20.19 + 24.942\times1.8018 = 65.130$; $3\ln(0.07/0.25) = -3.8189$।
\qty{6000}{bar} पर: $T = 54496/65.130 = \qty{836.7}{K} = \qty{563.6}{\celsius}$, फिर $P
= 22.80\times563.6 - 1093 + 8.314\times836.7\times(-3.8189)/6.21 = \qty{7479}{bar}$।
\qty{7479}{bar} पर: $T = \qty{570.6}{\celsius}$, $P = \qty{7603}{bar}$।
\qty{7603}{bar} पर: $T = \qty{571.2}{\celsius}$, $P = \qty{7614}{bar}$; एक बार और, $T =
\qty{571.3}{\celsius}$। चट्टान लगभग \qty{571}{\celsius} और
\qty{7.6}{kbar} पर साम्य में आई। (ख) \qty{571}{\celsius} पर कायनाइट $=$ सिलिमेनाइट रेखा $3.67
+ 0.71\times2.10 = \qty{5.16}{kbar}$ पर है: कायनाइट। (ग) दाबमापी
कायनाइट के लिए अंशांकित किया गया था; सिलिमेनाइट के साथ अभिक्रिया, और उसकी रेखा, भिन्न होती हैं, क्योंकि दोनों बहुरूपों की
गिब्स ऊर्जा समान नहीं है।
\end{solution}

\begin{solution}{exo:b2:pt-paths:11}
(क) $\tau = L^2/(\pi^2\kappa)$: \qty{35}{km} के लिए $1.225\times10^9/(9.870\times10^{-6}) =
\qty{1.24e14}{s} = \qty{3.9}{Myr}$; \qty{70}{km} के लिए चार गुना अधिक,
\qty{15.7}{Myr}। (ख) सत्तर मिलियन वर्ष इन समयों के चार से अठारह गुना हैं: ऊष्मा
मोटी हुई भूपर्पटी में से उससे कहीं तेज़ी से गुज़रती है जितनी तेज़ी से अपरदन उसे हटाता है। (ग) चट्टान
अपनी सबसे बड़ी गहराई पर एक साथ, चढ़ाव में ही, पहुँच जाती है; ऊष्मा अगले
दसियों मिलियन वर्षों में आती है, जबकि अपरदन उसे धीरे-धीरे ऊपर ला रहा होता है: विसंपीडन के दौरान
कुछ समय तक तापमान बढ़ता रहता है। (घ) सात मिलियन वर्षों में अपरदन
चट्टान को चालन द्वारा गर्म होने से पहले ही ऊपर ले आता: एक निम्नतर तापीय शिखर, जो
जल्दी, उच्चतर दाब पर, पहुँचा जाता, और सतह तक एक ठंडा वापसी-पथ।
\end{solution}

\begin{solution}{exo:b2:pt-paths:12}
(क) भूपर्पटी का भार $35000\times2800\times9.81 = \qty{0.961}{GPa}$ है; शेष
\qty{1.839}{GPa} के लिए $1.839\times10^9/(3300\times9.81) = \qty{56.8}{km}$ प्रावार चाहिए: गहराई
\qty{91.8}{km}, जैसे मूल अध्ययन के \qty{90}{km} या अधिक। (ख) $(4.0 -
0.961)\times10^9/(3300\times9.81) = \qty{93.9}{km}$, गहराई \qty{128.9}{km}।
(ग) $700/91.8 = \qty{7.6}{\celsius\per\km}$: दाब और तापमान का बहुत ऊँचा
अनुपात, क्षेपण क्षेत्रों की श्रेणी। (घ) $(91.8 - 35)/34 = \qty{1.7}{Myr}$।
\end{solution}

\begin{solution}{pb:b2:pt-paths:1}
\noindent\textbf{भाग I।}

\textbf{1.} गार्नेट क्रोड $0.080/0.760 = 0.1053$, बायोटाइट $0.40/0.58 = 0.6897$; $K_D =
0.1527$, $\ln K_D = -1.8793$।

\textbf{2.} किनारा $0.120/0.720 = 0.1667$, आधात्री बायोटाइट $0.46/0.53 = 0.8679$; $K_D =
0.1921$, $\ln K_D = -1.6497$।

\textbf{3.} $3R = 24.942$; अंश $52630 + 0.311\times6000 = 54496$; हर $20.19
+ 24.942\times1.8793 = 67.064$; $T = \qty{812.6}{K} = \qty{539.5}{\celsius}$।

\textbf{4.} हर $20.19 + 24.942\times1.6497 = 61.337$; $T = 54496/61.337 =
\qty{888.5}{K} = \qty{615.3}{\celsius}$।

\textbf{5.} $0.311\times1000/67.064 = \qty{4.6}{K}$ क्रोड के लिए और $0.311\times1000/61.337
= \qty{5.1}{K}$ किनारे के लिए।

\textbf{6.} तापमापी उन खनिजों पर लागू होता है जो एक साथ साम्य में थे।
क्रोड में अंतर्वेशित बायोटाइट क्रोड के उगते समय शेष चट्टान से अलग बंद हो गया;
आधात्री बायोटाइट ने अंतिम विनिमय किनारे के साथ किया। क्रोड को आधात्री के साथ जोड़ने से
दो क्षण मिल जाएँगे।

\noindent\textbf{भाग II।}

\textbf{7.} क्रोड $22.80\times539.5 - 1093 = \qty{11208}{bar}$; किनारा $22.80\times615.3 -
1093 = \qty{12936}{bar}$।

\textbf{8.} $3\ln(0.050/0.25) = -4.8283$; $RT = 8.314\times812.65 = \qty{6756.4}{J}$; $P =
11208 - 6756.4\times4.8283/6.21 = 11208 - 5253 = \qty{5955}{bar}$।

\textbf{9.} $3\ln(0.060/0.28) = -4.6213$; $RT = 8.314\times888.45 = \qty{7386.6}{J}$; $P
= 12936 - 5497 = \qty{7439}{bar}$।

\textbf{10.} \qty{7439}{bar} पर: $T = (52630 + 2314)/61.337 = \qty{895.8}{K} =
\qty{622.6}{\celsius}$; उस तापमान पर $P = 13102 - 5542 = \qty{7560}{bar}$; फिर $T
= \qty{623.2}{\celsius}$। किनारा लगभग \qty{623}{\celsius} और \qty{7.56}{kbar} पर उगा।

\textbf{11.} \qty{5954}{bar} पर, $T = \qty{539.2}{\celsius}$ और $P = \qty{5950}{bar}$:
क्रोड लगभग \qty{539}{\celsius} और \qty{5.95}{kbar} पर उगा। पहला अनुमान,
\qty{6}{kbar}, पहले ही निकट था, और तापमापी प्रति किलोबार केवल \qty{4.6}{K} खिसकता है।

\textbf{12.} क्रोड $5950\times10^5/(2800\times9.81) = \qty{21.7}{km}$; किनारा
$7560\times10^5/27468 = \qty{27.5}{km}$।

\textbf{13.} \qty{539}{\celsius} पर रेखा $3.67 + 0.39\times2.10 =
\qty{4.49}{kbar}$ पर है, \qty{5.95}{kbar} से नीचे: कायनाइट, जो क्रोड में अंतर्वेशित कायनाइट से
संगत है। \qty{623}{\celsius} पर वह $5.77 + 0.23\times2.09 = \qty{6.25}{kbar}$ पर है,
\qty{7.56}{kbar} से नीचे: फिर कायनाइट।

\noindent\textbf{भाग III।}

\textbf{14.} $+\qty{84}{\celsius}$ और $+\qty{1.6}{kbar}$ (लगभग \qty{6}{km} अधिक गहरा):
चट्टान दब रही थी और गर्म हो रही थी, अपने \omterm{def:b2:pt-paths:paths}{प्रगामी पथ} पर।

\textbf{15.} \qty{623}{\celsius} या अधिक पर, \qty{6.25}{kbar} से नीचे: किनारा उगने के बाद
चट्टान ने कम से कम \qty{1.3}{kbar} खोया, यानी लगभग \qty{5}{km} ऊपरी भार, जबकि वह कम से कम
उतनी ही गर्म रही। उसका \omterm{def:b2:pt-paths:exhumation}{उत्खनन} तब हुआ जब वह अभी गर्म थी।

\textbf{16.} क्लोराइट एक निम्न-श्रेणी, जलयोजित खनिज है: वह बाद के \omterm{def:b2:pt-paths:paths}{प्रतिगामी
पथ} को अंकित करता है, उथली गहराई पर, जल के साथ, ठंडा होना।

\textbf{17.} किनारे तक दबना और गर्म होना, ठंडा हुए बिना विसंपीडन, फिर ठंडा होना:
तापमान को दाईं ओर और दाब को ऊपर रखकर खींचने पर पथ दक्षिणावर्त मुड़ता है।
दाब-शिखर पहले आया, किनारे पर; तापीय शिखर उसी समय या बाद में, विसंपीडन के
दौरान।

\textbf{18.} हाँ: मॉडल भी पहले अपने दाब-शिखर पर पहुँचता है, और आरोहण के दौरान अपने तापीय शिखर पर,
यानी \qty{593}{\celsius}, \qty{7.8}{kbar} पर, जो शिस्ट के किनारे की
दशाओं के निकट है। मोटा होने के बाद धीमा अपरदन शिस्ट के पथ की व्याख्या करता है।

\noindent\textbf{भाग IV।}

\textbf{19.} $(623 - 405)/(412.0 - 398.0) = 218/14.0 = \qty{15.6}{\celsius}$ प्रति मिलियन
वर्ष।

\textbf{20.} चट्टान मान्यता से लगभग डेढ़ गुना तेज़ी से ठंडी हुई। जो खनिज
तेज़ी से ठंडा होता है उसके पास अपना आर्गन खोने के लिए कम समय होता है, इसलिए वह उच्चतर तापमान पर बंद होता है:
वास्तविक अवरोध-तापमान \qty{405}{\celsius} से कुछ ऊपर है, और बंद होना
कुछ अधिक गहराई पर।

\textbf{21.} $(405 - 10)/30 = \qty{13.2}{km}$।

\textbf{22.} $27.5 - 13.2 = \qty{14.3}{km}$ (अपूर्णांकित गहराइयों, 27.52 और 13.17, के साथ
\qty{14.35}{km})।

\textbf{23.} $\sqrt{1.5^2 + 1.0^2} = \qty{1.8}{Myr}$: $14.0 \pm \qty{1.8}{Myr}$, अतः दर इन दोनों के बीच है:
$14.35/15.8 = 0.91$ और $14.35/12.2 = \qty{1.18}{km/Myr}$।

\textbf{24.} $14.35/14.0 = \qty{1.03}{km}$ प्रति मिलियन वर्ष, यानी
\textbf{\qty[text-series-to-math=true]{1.0}{mm.yr^{-1}}}, जो आल्प्स की सबसे गहरी चट्टानों के \omterm{def:b2:pt-paths:exhumation}{उत्खनन} के
पहले चरण से लगभग तीस गुना धीमा है, और मोटा होने वाले मॉडल की अपरदन-दर का
लगभग दोगुना।
\end{solution}
